\section{Introduction}

India's police and criminal-justice information systems already operate at national scale. The Crime and Criminal Tracking Network and Systems (CCTNS) provides a digital backbone for police-station processes, and an official Press Information Bureau release reports that all 17,798 police stations were using CCTNS as of February 1, 2026 \cite{pib_cctns_2026}. The same source describes CCTNS support for FIRs, chargesheets, state-hosted applications, near-real-time replication to a National Data Centre, crime and criminal search, standardized Integrated Investigation Forms, and master codes for states, districts, police stations, acts, and sections \cite{pib_cctns_2026}. The Inter-Operable Criminal Justice System (ICJS) further connects the police/CCTNS pillar with courts, prisons, forensics, and prosecution through a Data Sharing Matrix \cite{mha_icjs_2026}. Therefore, this paper does not treat Indian policing as a blank technical space. It starts from the position that CCTNS and ICJS-style infrastructure already exist and should be complemented, not replaced.

The problem addressed in this paper is controlled access to sensitive inter-agency police and criminal-justice records. These records may include FIR details, witness and victim information, juvenile records, forensic reports, cybercrime complaints, case diary material, evidence references, and court or prosecution-linked information. Such records are different from ordinary administrative data because access may affect privacy, investigation integrity, victim protection, and accountability. A broad role such as ``police officer'' or ``agency user'' is not enough to decide whether a record should be released. The decision may depend on the requester's role, rank, station, jurisdiction, case assignment, purpose of access, record sensitivity, credential status, approval state, emergency context, time window, and connection to court, prosecution, or forensic workflow.

\textbf{Problem statement:} Given an inter-agency request for a sensitive police or criminal-justice record, the system must decide whether to allow, deny, or escalate the request using contextual policy rules, while preserving an auditable and explainable record of the decision. The record should show what was requested, which policy version was applied, which attributes affected the decision, what explanation was generated, and which approval or review event was associated with the request. This must be done without placing raw sensitive records on a public or shared ledger.

This problem requires three equal technical pillars. The first pillar is security and access control. Static role-based access control (RBAC) is useful, but it is too coarse for many sensitive inter-agency workflows. Attribute-Based Access Control (ABAC) evaluates subject, object, action, and environmental attributes against policy rules or relationships \cite{nist_abac_2014}. In this work, ABAC is combined with policy-based access control (PBAC) ideas such as policy versioning, purpose constraints, credential revocation, approval requirements, and superior review. This keeps the authorization decision in a policy layer rather than shifting trust to the audit layer.

The second pillar is blockchain-based auditability. Blockchain is not used here to store raw police records or to provide privacy by itself. Instead, it is used as a tamper-evident audit layer for known agencies. NISTIR 8403 describes blockchain-based access-control systems in terms of decentralization, high confidence, and tamper resistance, while also noting practical implementation considerations \cite{nist_blockchain_access_2022}. Hyperledger Fabric is a relevant reference point because it supports permissioned blockchain networks among known organizations \cite{androulaki_fabric_2018}. In the proposed design, raw records remain off-chain under agency control, while the ledger stores minimized commitments such as request identifiers, policy identifiers, approval events, model or rule versions, explanation hashes, and record commitments.

The third pillar is explainable AI (XAI). Access decisions in public-safety workflows must be understandable to officers, supervisors, auditors, and possibly court or prosecution stakeholders. In this paper, XAI is not treated as a decorative dashboard. It is treated as a logged audit artifact. Each decision should preserve the decision label, reason code, decisive attributes, policy version, model or rule version, counterfactual information where applicable, and hashes that bind the explanation to the audit event. This position follows the high-stakes AI argument that interpretable models are preferable where feasible \cite{rudin_2019}, and recent law-enforcement XAI work that emphasizes stakeholder needs, AI literacy, and automation-bias risk \cite{zocholl_xai_law_enforcement_2025}.

This paper is intentionally not framed as individual crime prediction. Predictive policing research has shown that feedback loops can occur when observed or reported crime data is shaped by earlier policing activity \cite{ensign_feedback_2018}. Public NCRB Crime in India data is useful as aggregate reported/registered crime context, but it is not a public individual-level CCTNS/FIR access-decision dataset \cite{ncrb_2023}. Therefore, the technical focus of this work is secure, explainable, and auditable access governance, not predicting criminals, suspects, or future crimes.

Existing research covers important parts of the problem, but not the full workflow considered here. Blockchain has been studied for digital crime evidence management and lawful evidence chain-of-custody workflows \cite{kim_two_level_2021,li_lechain_2021}. Fabric-based ABAC and encrypted off-chain data sharing have been studied in generic data-sharing contexts \cite{zhao_fabric_abac_2022}. XAI and fairness literature has raised strong cautions for high-stakes criminal-justice systems \cite{rudin_2019,ensign_feedback_2018,zocholl_xai_law_enforcement_2025}. The remaining gap is a CCTNS/ICJS-compatible access-governance overlay that jointly studies contextual authorization, superior-review style escalation, off-chain sensitive-record commitments, permissioned blockchain-style audit, XAI artifact logging, adversarial audit attacks, and explanation reviewability.

To address this gap, this paper proposes SEBA-XAI, a Secure, Explainable, Blockchain-Audited access-governance overlay for sensitive inter-agency police record sharing. SEBA-XAI is designed to sit above existing agency systems. The proposed methodology uses synthetic CCTNS/ICJS-style access requests, a deterministic policy oracle, a permissioned blockchain-style audit layer, off-chain record commitments, and XAI artifacts linked to audit events. This synthetic setup is used because actual police access logs and sensitive records are not publicly available and should not be assumed for a first-stage academic prototype.

The initial contributions of this paper are as follows:
\begin{itemize}
    \item It formulates a CCTNS/ICJS-compatible access-governance problem for sensitive inter-agency police and criminal-justice records.
    \item It proposes the SEBA-XAI architecture combining RBAC/ABAC/PBAC-style policy enforcement, off-chain sensitive records, permissioned blockchain-style audit commitments, and logged XAI artifacts.
    \item It defines a reproducible synthetic access-request benchmark for allow, deny, and escalate decisions under controlled policy and attack conditions.
    \item It proposes evaluation measures for audit reconstruction, tamper detection, false allow and false deny behavior, explanation completeness, explanation stability, latency, storage overhead, and metadata exposure.
\end{itemize}

The scope of the paper is limited. It does not claim live deployment inside CCTNS or ICJS, does not use actual police records, does not store raw records on-chain, does not prove legal compliance, and does not claim state-of-the-art crime prediction. The aim is to write and evaluate a careful research architecture that can later be refined through supervisor feedback, stronger experiments, and domain validation.
